本实验研究一个动态特性未知且具有明显非线性的温度控制对象。实验首先利用阶跃响应建立线性近似模型，再按照给定的超调量和调节时间设计代数补偿控制器，最后分别通过 Simulink 和 dSpace 实验台评价闭环性能。除设定值跟踪外，还需要讨论入口温度 T1、空气质量流量和执行器限幅等对控制品质的影响[1]。
本报告严格按照 MRT 通用要求组织，包含实验动机、装置、过程、结果分析和误差讨论，并逐项回答实验任务 a) 至 i)[2]。封面上的姓名、学号、组号和日期保留为空白。
